\documentclass[10pt,a4paper]{amsart}
\usepackage[margin=19mm]{geometry}
\usepackage{amsmath,amssymb,mathtools}
\usepackage[scr=rsfs,cal=euler]{mathalfa}
\usepackage{xfrac}
\usepackage[unicode,hidelinks,bookmarks=false]{hyperref}
\newcommand{\ie}{i.e.\ }
\newcommand{\cf}{cf.\ }
\newcommand{\resp}{respectively\ }
\newcommand{\Subsection}{\subsection}
\providecommand{\nobreakdash}{\nobreak\hbox{-}}
\setlength{\emergencystretch}{2em}
\multlinegap0pt
\usepackage{amssymb}
\usepackage{mathtools}
\usepackage{tikz}
\usepackage[normalem]{ulem}
\newtheorem{proposition}{Proposition}
\newcommand{\C}{{\mathbb C}}
\DeclareMathOperator{\SL}{SL}
\title{A Lie group analog for the Monster Lie algebra}
\author{Lisa Carbone, Elizabeth Jurisich, Scott H. Murray}
\date{Source: arXiv:2311.11078v4 (CC BY 4.0)}
\begin{document}
\maketitle

\noindent\emph{Source and licence.} A Lie group analog for the Monster Lie algebra. Version arXiv:2311.11078v4 (CC BY 4.0), distributed by arXiv under the Creative Commons Attribution 4.0 International licence, http://creativecommons.org/licenses/by/4.0/, as recorded in the arXiv licence metadata for this exact version and shown on the abstract page https://arxiv.org/abs/2311.11078v4. Full licence notice: \texttt{licenses/arxiv-2311.11078-CC-BY-4.0.txt}.\par\smallskip

\noindent\textit{Editorial scope.} Counted unit: the article's proposition nontriv (source0.tex lines 3201--3203). The article's complete original proof of it is delivered with it (source0.tex lines 3205--3344, one \texttt{proof} environment of 140 lines). No mathematics is written, repaired or re-derived: the statement and the proof are the article's own lines, in the article's own notation, and no retained line is substituted or re-flowed.  No further source-local context is needed: every cross-reference of every retained line resolves inside the two delivered spans.  Nothing was omitted from any retained span, so the package carries no omission record.  No retained line contains a citation, so the wrapper carries no bibliography and no bibliography entry.  No retained line contains a figure environment, an image-inclusion command or drawing code, so no image is delivered and the package declares no asset dependency.  Editorial formatting only, in addition: an AMS \texttt{amsart} preamble that restates the article's own declarations and macros used by the retained lines, namely the environments \texttt{proposition} on separate counters, together with the standard packages those lines need.  The retained environments print in this document as Proposition 1 in order of appearance, whereas the article prints the same items with the numbering of its own full document.  The complete native source of the article as distributed in the arXiv e-print for this exact version is bundled unchanged as \texttt{sources/150175/source0.tex}.
 {\color{blue} \begin{proposition}\label{nontriv}
The group $G = G_{0,2k;\,1,2k}$ is nontrivial.
\end{proposition}

\begin{proof} We construct  injective group homomorphisms,   from a pair of subgroups of $ G_{0,2k;\,1,2k}$ to $\SL_3(\mathbb{C})$ as follows. 

For any $u, v \in \mathbb{C}$ and $s \in \mathbb{C}^{\times}$ we set
$$\rho_1(h_1)=\begin{pmatrix} 
s & 0 & 0 \\ 
0 & 1 & 0 \\ 
0 & 0 & s^{-1} \end{pmatrix}, 
\rho_1(X_{-1}(u)) = \begin{pmatrix} 1 & u & 0 \\ 0 & 1 & 0 \\ 0 & 0 & 1 \end{pmatrix}, $$
$$
\rho_1(X_{1,2k}(u)) = \begin{pmatrix} 1 & 0 & u \\ 0 & 1 & 0 \\ 0 & 0 & 1 \end{pmatrix}, \rho_1(Y_{0,2k}(u)) = \begin{pmatrix} 1 & 0 & 0 \\ 0 & 1 & 0 \\ 0 & u & 1 \end{pmatrix}.$$

We simultaneously upper-triangularize these  matrices as follows. We choose the generators of the one parameter subgroups:
$E_{12}$ for $\rho(X_{-1}(u))$
$E_{13}$ for $\rho(X_{1,2k}(u))$ and 
$E_{32}$ for $\rho(Y_{0,2k}(u))$.


Note that $[E_{32}, E_{13}] = -E_{12}$, and $E_{12}$ commutes with both $E_{32}$ and $E_{13}$. Thus, $\{E_{12}, E_{13}, E_{32}\}$ spans a nilpotent Lie subalgebra isomorphic to the 3-dimensional Heisenberg algebra $\mathfrak{h}_3$. 


Let $\{e_1, e_2, e_3\}$ be the standard basis for $\mathfrak{sl}_2(\C)$. We define a new basis $\{v_1, v_2, v_3\}$ as follows:
\[
v_1 = e_1, \quad v_2 = e_3, \quad v_3 = -e_2
\]


The corresponding change-of-basis matrix $P$  and its inverse $P^{-1}$ are:
\[
P_1 = \begin{pmatrix} 1 & 0 & 0 \\ 0 & 0 & -1 \\ 0 & 1 & 0 \end{pmatrix}, \quad P_1^{-1} = \begin{pmatrix} 1 & 0 & 0 \\ 0 & 0 & 1 \\ 0 & -1 & 0 \end{pmatrix}
\]



{ For $\rho_1(X_{-1}(u))$} we have
\[
P_1^{-1} \begin{pmatrix} 1 & u & 0 \\ 0 & 1 & 0 \\ 0 & 0 & 1 \end{pmatrix} P_1 = \begin{pmatrix} 1 & 0 & -u \\ 0 & 1 & 0 \\ 0 & 0 & 1 \end{pmatrix}.
\]

{For $\rho_1(X_{1,2k}(u))$} we have
\[
P_1^{-1} \begin{pmatrix} 1 & 0 & u \\ 0 & 1 & 0 \\ 0 & 0 & 1 \end{pmatrix} P_1 = \begin{pmatrix} 1 & u & 0 \\ 0 & 1 & 0 \\ 0 & 0 & 1 \end{pmatrix}.
\]

{For $\rho_1(Y_{0,2k}(u))$} we have
\[
P_1^{-1} \begin{pmatrix} 1 & 0 & 0 \\ 0 & 1 & 0 \\ 0 & u & 1 \end{pmatrix} P_1 = \begin{pmatrix} 1 & 0 & 0 \\ 0 & 1 & -u \\ 0 & 0 & 1 \end{pmatrix}.
\]

In this new basis, all matrices are elements of the standard unipotent subgroup of $\mathrm{SL}_3(\mathbb{C})$.


These generate a copy $B_1$ of the Borel subgroup of $\SL_3(\C)$.



We define
\begin{align*}
\rho_2(h_2) &= \begin{pmatrix}1&0&0\\ 0&s&0\\ 0&0&s^{-1}\end{pmatrix}, &
\rho_2(Y_{-1}(u)) &= \begin{pmatrix}1&0&0\\ u&1&0\\ 0&0&1\end{pmatrix}, \\
\rho_2(X_{0,2k}(u)) &= \begin{pmatrix}1&0&0\\ 0&1&-u\\ 0&0&1\end{pmatrix}, &
\rho_2(Y_{1,2k}(u)) &= \begin{pmatrix}1&0&0\\ 0&1&0\\ u&0&1\end{pmatrix}.
\end{align*}

Let $\{e_1, e_2, e_3\}$ be the standard basis for $\mathbb{C}^3$. We define a new basis $\{v_1, v_2, v_3\}$ where $v_1=e_2$, $v_2=e_3$, and $v_3=e_1$. The corresponding change-of-basis matrix $P_2$ and its inverse $P_2^{-1}$ are:
\[
P_2 = \begin{pmatrix} 0 & 0 & 1 \\ 1 & 0 & 0 \\ 0 & 1 & 0 \end{pmatrix}, \qquad
P_2^{-1} = \begin{pmatrix} 0 & 1 & 0 \\ 0 & 0 & 1 \\ 1 & 0 & 0 \end{pmatrix}.
\]

We similarly upper-triangularize these matrices using $P_2$ to obtain:

For $\rho_2(Y_{-1}(u))$ we have:
\[
P_2^{-1}\begin{pmatrix}1&0&0\\ u&1&0\\ 0&0&1\end{pmatrix}P_2=\begin{pmatrix}1&0&u\\ 0&1&0\\ 0&0&1\end{pmatrix}.
\]
For $\rho_2(X_{0,2k}(u))$ we have:
\[
P_2^{-1}\begin{pmatrix}1&0&0\\ 0&1&-u\\ 0&0&1\end{pmatrix}P_2=\begin{pmatrix}1&-u&0\\ 0&1&0\\ 0&0&1\end{pmatrix}.
\]
For $\rho_2(Y_{1,2k}(u))$ we have:
\[
P_2^{-1}\begin{pmatrix}1&0&0\\ 0&1&0\\ u&0&1\end{pmatrix}P_2=\begin{pmatrix}1&0&0\\ 0&1&u\\ 0&0&1\end{pmatrix}.
\]
These generate $B_2$, a copy of the standard Borel subgroup of $SL_3(\mathbb{C})$.


\newpage


Now  we form the generalized amalgamated product $B$ of $B_1,B_2, H, G_{0,2k}, G_{1,2k}, G_{-1}$ over their common intersections.








{\color{red} We set
$$A=\langle \widetilde{w}_{-1}, h_1, h_2 \rangle$$
and
$$H=\langle h_1,h_2\rangle.$$
Conjugation by $\widetilde{w}_{-1}$ normalizes $H$  and $\widetilde{w}_{-1}^2 \in H$. Thus, $H \lhd A$, and  $A/H \cong W$.

 We also form $G'_{0,2k:1,2k}=A*_H B$.
 
Our strategy for showing that $G_{0,2k:1,2k}$ is nontrivial is to show that $G_{0,2k:1,2k}=G'_{0,2k:1,2k}$. Then since $G'_{0,2k:1,2k}=A*_H B$ is an amalgam, it will follow that $G_{0,2k:1,2k}$ is nontrivial.


Here is the rigorous mathematical breakdown of why $G_{0,2k;1,2k}$ is actually a proper quotient of the amalgam $A *_H B$.

By definition, the base group $B$ is the generalized amalgamated product of $B_1, B_2, H, G_{0,2k}, G_{1,2k}$, and $G_{-1}$. Because $B$ explicitly contains $G_{-1}$ as a subgroup, $B$ must inherently satisfy all relations defining $G_{-1}$. One of those defining relations (Re:0) strictly defines the Weyl element via the unipotent generators:
$$\widetilde{w}_{-1} = X_{-1}(1)Y_{-1}(-1)X_{-1}(1)$$

Therefore, the element $\widetilde{w}_{-1}$ already exists strictly inside $B$. Because the group $A$ is defined as $A = \langle \widetilde{w}_{-1}, H \rangle$, and both $\widetilde{w}_{-1}$ and $H$ are contained in $B$, it follows that $A$ is effectively a subgroup of $B$. If we attempt to form the abstract free product with amalgamation $G' = A *_H B$, the operation treats the factor $A$ and the factor $B$ as distinct entities outside of $H$. This means $G'$ will contain two independent copies of the Weyl element:

$w_A \in A \setminus H$ and $w_B = X_{-1}(1)Y_{-1}(-1)X_{-1}(1) \in B \setminus H$

In the amalgam $A *_H B$, these two elements are completely distinct. This shatters the fundamental relation $\widetilde{w}_{-1} = X_{-1}(1)Y_{-1}(-1)X_{-1}(1)$ required by $G_{0,2k;1,2k}$.

Even if we collapse $A$ into $B$ (meaning $A *_H B$ simply defaults to $B$), the base group $B$ still does not equal $G_{0,2k;1,2k}$. The presentation for $G_{0,2k;1,2k}$ includes critical cross-generator relations where the Weyl element acts as an automorphism on the imaginary root unipotents. 

For example:
$$\widetilde{w}_{-1} X_{0,2k}(u) \widetilde{w}_{-1}^{-1} = X_{1,2k}(u)$$In a generalized amalgamated product, relations only hold if they exist within one of the vertex subgroups (factors) or are explicitly forced by the amalgamation over the edge subgroups (intersections).$\widetilde{w}_{-1}$ lives entirely within the $G_{-1}$ factor.$X_{0,2k}(u)$ lives in the $G_{0,2k}$ and $B_2$ factors.$X_{1,2k}(u)$ lives in the $G_{1,2k}$ and $B_1$ factors.There is no single factor in the construction of $B$ that contains $\widetilde{w}_{-1}$, $X_{0,2k}$, $X_{1,2k}$, and the conjugation relations between them. 

In the amalgam $B$, the sequence $\widetilde{w}_{-1} X_{0,2k}(u) \widetilde{w}_{-1}^{-1}$ is simply an irreducible word in normal form. It does not algebraically evaluate to $X_{1,2k}(u)$.

Because the amalgam $G'_{0,2k;1,2k} = A *_H B$ lacks the cross-factor conjugation relations and duplicates the Weyl element, $G_{0,2k;1,2k}$ is a proper quotient of $G'_{0,2k;1,2k}$.Consequently, showing that the amalgam $G'$ is nontrivial does not mathematically establish that the quotient group $G$ is nontrivial. 

}









\end{proof}

\end{document}
